\documentclass[10pt,a4paper]{amsart}
\usepackage{amsmath,amssymb,amsthm,graphicx}
% For commutative diagrams we use Xy-pic
\usepackage[all]{xy}

% We use 2cell for 2-commutative diagrams.
\xyoption{2cell}
% Theorem environments.
%
\theoremstyle{plain}
\newtheorem{theorem}[subsection]{Theorem}
\newtheorem{proposition}[subsection]{Proposition}
\newtheorem{lemma}[subsection]{Lemma}

\theoremstyle{definition}
\newtheorem{definition}[subsection]{Definition}
\newtheorem{example}[subsection]{Example}
\newtheorem{exercise}[subsection]{Exercise}
\newtheorem{situation}[subsection]{Situation}

\theoremstyle{remark}
\newtheorem{remark}[subsection]{Remark}
\newtheorem{remarks}[subsection]{Remarks}

\numberwithin{equation}{subsection}

% Macros
%
\def\lim{\mathop{\mathrm{lim}}\nolimits}
\def\colim{\mathop{\mathrm{colim}}\nolimits}
\def\Spec{\mathop{\mathrm{Spec}}}
\def\Hom{\mathop{\mathrm{Hom}}\nolimits}
\def\Ext{\mathop{\mathrm{Ext}}\nolimits}
\def\SheafHom{\mathop{\mathcal{H}\!\mathit{om}}\nolimits}
\def\SheafExt{\mathop{\mathcal{E}\!\mathit{xt}}\nolimits}
\def\Sch{\mathit{Sch}}
\def\Mor{\mathop{\mathrm{Mor}}\nolimits}
\def\Ob{\mathop{\mathrm{Ob}}\nolimits}
\def\Sh{\mathop{\mathit{Sh}}\nolimits}
\def\NL{\mathop{N\!L}\nolimits}
\def\CH{\mathop{\mathrm{CH}}\nolimits}
\def\proetale{{pro\text{-}\acute{e}tale}}
\def\etale{{\acute{e}tale}}
\def\QCoh{\mathit{QCoh}}
\def\Ker{\mathop{\mathrm{Ker}}}
\def\Im{\mathop{\mathrm{Im}}}
\def\Coker{\mathop{\mathrm{Coker}}}
\def\Coim{\mathop{\mathrm{Coim}}}

% Boxtimes
%
\DeclareMathSymbol{\boxtimes}{\mathbin}{AMSa}{"02}

%
% Macros for moduli stacks/spaces
%
\def\QCohstack{\mathcal{QC}\!\mathit{oh}}
\def\Cohstack{\mathcal{C}\!\mathit{oh}}
\def\Spacesstack{\mathcal{S}\!\mathit{paces}}
\def\Quotfunctor{\mathrm{Quot}}
\def\Hilbfunctor{\mathrm{Hilb}}
\def\Curvesstack{\mathcal{C}\!\mathit{urves}}
\def\Polarizedstack{\mathcal{P}\!\mathit{olarized}}
\def\Complexesstack{\mathcal{C}\!\mathit{omplexes}}
% \Pic is the operator that assigns to X its picard group, usage \Pic(X)
% \Picardstack_{X/B} denotes the Picard stack of X over B
% \Picardfunctor_{X/B} denotes the Picard functor of X over B
\def\Pic{\mathop{\mathrm{Pic}}\nolimits}
\def\Picardstack{\mathcal{P}\!\mathit{ic}}
\def\Picardfunctor{\mathrm{Pic}}
\def\Deformationcategory{\mathcal{D}\!\mathit{ef}}
\usepackage{cleveref}
\crefname{originaltheorem}{original theorem}{original theorems}
\Crefname{originaltheorem}{Original theorem}{Original theorems}
\makeatletter
\newlabel{algebra-theorem-lazard}{{058G}{source}{}{sourcepointer-1}{}}
\newlabel{algebra-theorem-lazard@cref}{{[originaltheorem][0][]058G}{[1][0][]source}{}{}{}}
\makeatother
\crefname{originallemma}{original lemma}{original lemmas}
\Crefname{originallemma}{Original lemma}{Original lemmas}
\makeatletter
\newlabel{more-algebra-lemma-last-one-flat}{{0653}{source}{}{sourcepointer-2}{}}
\newlabel{more-algebra-lemma-last-one-flat@cref}{{[originallemma][0][]0653}{[1][0][]source}{}{}{}}
\makeatother
\crefname{originalreference}{original reference}{original references}
\Crefname{originalreference}{Original reference}{Original references}
\makeatletter
\newlabel{theorem-colimit}{{09YQ}{source}{}{sourcepointer-3}{}}
\newlabel{theorem-colimit@cref}{{[originalreference][0][]09YQ}{[1][0][]source}{}{}{}}
\makeatother
\title{Global sections of a derived tensor product with a constant complex of abelian groups pull the constant factor out}
\author{The Stacks Project Authors}
\date{stacks-project:etale-cohomology.tex:local-snapshot-20261011; GNU FDL 1.2 or later}
\begin{document}
\maketitle
\noindent\textit{Editorial note.} All delivered mathematical text is the author's own, copied from the pinned Stacks Project chapter file etale-cohomology.tex; no statement or proof step was rewritten, reordered or removed.\par\medskip
\noindent\textit{Editorial note.} The counted proof is the author's complete argument: the proof is written out in full: the choice of the torsion free representative justified through global dimension one, the truncation reductions using the two displayed vanishing statements, and the final reduction to a finite free module by Lazard's theorem and the colimit theorem.\par\medskip
\noindent\textit{Editorial note.} The label lemma-pull-out-constant names the counted statement, and the proof environment immediately following it in the source is delivered verbatim as its proof.\par\medskip
\noindent\textit{Editorial note.} The counted statement is carried by the author's own lemma environment in the chapter source.\par\medskip
\noindent\textit{Editorial note.} The statement and proof are copied from the pinned chapter file; the three cited inputs are kept as external pointers with their canonical Stacks tags.\par\medskip
\noindent\textit{Editorial note.} The preamble is the chapter's own native preamble excerpt already used for this source: the Xy-pic setup (native preamble lines 14-18) and the native theorem declarations with the mathematical macros (native preamble lines 154-217).\par\medskip
\noindent\textit{Allowed standard prerequisites (editorial).}\begin{itemize}
\item Constant sheaf underline K associated to a complex of abelian groups. Same chapter, the constant sheaf on X\_etale associated to a complex K of abelian groups and its derived tensor product with an object of D(X\_etale). Supplies the object E tensor\textasciicircum{}L\_Z underline K whose global sections are computed.
\item K-flat resolutions over Z and the last-one-flat criterion. More on Algebra, the existence of K-flat resolutions and the flatness criterion used to replace K by a bounded below complex of torsion free modules. Supplies the representative K\textasciicircum{}bullet = tau\_\{>= b-1\}L\textasciicircum{}bullet of K by torsion free modules, using that Z has global dimension one.
\item Lazard's theorem and the colimit theorem for etale cohomology. Algebra, the description of a flat module as a filtered colimit of finite free modules, and the same chapter's theorem computing cohomology of filtered colimits. Supplies the reduction from a single flat module to a finite free module in the last step of the proof.
\end{itemize}\medskip
\ifdefined\hypertarget\hypertarget{sourcepointer-1}{}\fi
\noindent\textit{Source reference.} Original theorem 058G: sources/stacks-chow/tags.txt, source lines 6779--6779. Prerequisite pointer invoked at proof line(s) 18995; it names a standard fact of a sibling chapter that the counted argument uses, and every step of the counted proof built on it is retained verbatim in the delivered body.\par\medskip
\ifdefined\hypertarget\hypertarget{sourcepointer-2}{}\fi
\noindent\textit{Source reference.} Original lemma 0653: sources/stacks-chow/tags.txt, source lines 7887--7887. Prerequisite pointer invoked at proof line(s) 18955; it names a standard fact of a sibling chapter that the counted argument uses, and every step of the counted proof built on it is retained verbatim in the delivered body.\par\medskip
\ifdefined\hypertarget\hypertarget{sourcepointer-3}{}\fi
\noindent\textit{Source reference.} Original reference 09YQ: sources/stacks-etale-cohomology/pinned-original-source.tex, source lines 7400--7400. Prerequisite pointer invoked at proof line(s) 18996; it names a standard fact of the same chapter that the counted argument uses, and every step of the counted proof built on it is retained verbatim in the delivered body.\par\medskip
\begin{lemma}
\label{lemma-pull-out-constant}
Let $X$ be a quasi-compact and quasi-separated scheme.
Let $E \in D^+(X_\etale)$ and $K \in D^+(\mathbf{Z})$.
Then
$$
R\Gamma(X, E \otimes_\mathbf{Z}^\mathbf{L} \underline{K}) =
R\Gamma(X, E) \otimes_\mathbf{Z}^\mathbf{L} K
$$
\end{lemma}

\begin{proof}
Say $H^i(E) = 0$ for $i \geq a$ and $H^j(K) = 0$ for $j \geq b$.
We may represent $K$ by a bounded below complex
$K^\bullet$ of torsion free $\mathbf{Z}$-modules.
(Choose a K-flat complex $L^\bullet$ representing $K$ and then
take $K^\bullet = \tau_{\geq b - 1}L^\bullet$. This works because
$\mathbf{Z}$ has global dimension $1$. See
More on Algebra, Lemma \ref{more-algebra-lemma-last-one-flat}.)
We may represent $E$ by a bounded below complex
$\mathcal{E}^\bullet$.
Then $E \otimes_\mathbf{Z}^\mathbf{L} \underline{K}$ is
represented by
$$
\text{Tot}(\mathcal{E}^\bullet \otimes_\mathbf{Z} \underline{K}^\bullet)
$$
Using distinguished triangles
$$
\sigma_{\geq -b + n + 1}K^\bullet
\to K^\bullet \to
\sigma_{\leq -b + n}K^\bullet
$$
and the trivial vanishing
$$
H^n(X,
\text{Tot}(\mathcal{E}^\bullet \otimes_\mathbf{Z}
\sigma_{\geq -a + n + 1}\underline{K}^\bullet) = 0
$$
and
$$
H^n(R\Gamma(X, E)
\otimes_\mathbf{Z}^\mathbf{L} \sigma_{\geq -a + n + 1}K^\bullet) = 0
$$
we reduce to the case where $K^\bullet$ is a bounded complex of
flat $\mathbf{Z}$-modules. Repeating the argument we reduce to the
case where $K^\bullet$ is equal to a single flat $\mathbf{Z}$-module
sitting in some degree. Next, using the stupid
trunctions for $\mathcal{E}^\bullet$
we reduce in exactly the same manner to the case where
$\mathcal{E}^\bullet$ is a single abelian sheaf sitting
in some degree. Thus it suffices to show that
$$
H^n(X, \mathcal{E} \otimes_\mathbf{Z} \underline{M}) =
H^n(X, \mathcal{E}) \otimes_\mathbf{Z} M
$$
when $M$ is a flat $\mathbf{Z}$-module and $\mathcal{E}$
is an abelian sheaf on $X$. In this case we
write $M$ is a filtered colimit of finite free $\mathbf{Z}$-modules
(Lazard's theorem, see Algebra, Theorem \ref{algebra-theorem-lazard}).
By Theorem \ref{theorem-colimit} this reduces us to the case of finite free
$\mathbf{Z}$-module $M$ in which case the result is trivially true.
\end{proof}

\end{document}
